% -----------------------------------------------------------------------------
% This file is automatically generated.
% Do not edit manually.
% Source: notebooks/support/regression-analysis/support.Rmd
% -----------------------------------------------------------------------------


\noindent The Breusch-Godfrey test is executed through \ttblue{bgtest}.

\par\addvspace{\topsep}
\begingroup
\fvset{listparameters={%
  \setlength{\topsep}{0pt}%
  \setlength{\partopsep}{0pt}%
  \setlength{\parsep}{0pt}%
  \setlength{\itemsep}{0pt}%
}}
\begin{Verbatim}[breaklines=true,formatcom=\color{ada_blue}]
bg_lm, bg_p, _, _ = acorr_breusch_godfrey(ussteamco_mod_3, nlags=3)
print({'LM statistic': bg_lm, 'p_value': bg_p})
\end{Verbatim}
\begin{Verbatim}[breaklines=true,formatcom=\color{ada_light_blue}]

	Breusch-Godfrey test for serial correlation of order up
	to 3

data:  ussteamco.mod.3
LM test = 7.1466, df = 3, p-value = 0.06737
\end{Verbatim}
\endgroup
\par\addvspace{\topsep}
\noindent The $p$ value of the test statistic is 0.06737. Therefore, we reject the null hypothesis at levels below 0.06737, suggesting marginal evidence of autocorrelation in the residuals.
